%!TEX root = ../main.tex
% Declaración sobre el uso de inteligencia artificial.
% Exigencia de la cátedra (clase 5): el uso de IA está permitido y
% fomentado, pero debe declararse. Esta página resume el uso; el
% registro completo de prompts va en el apéndice D.

\begin{guia}[Guía de la cátedra: política de uso de IA]
La IA es parte del estado del arte de la industria y su uso está
permitido, con condiciones:
\begin{enumerate}
    \item \textbf{Declararlo}: acá y en la Metodología (qué tareas fueron
          asistidas y con qué herramientas).
    \item \textbf{Autoría}: las ideas, decisiones y conclusiones son del
          estudiante. La IA asiste; no decide.
    \item \textbf{Verificación}: toda cifra, cita o referencia sugerida por
          una IA se confirma en la fuente original antes de usarla.
    \item \textbf{Trazabilidad}: los prompts relevantes se guardan en una
          carpeta compartida con la cátedra.
    \item \textbf{Comprensión}: no se incorpora texto que no se haya leído
          y entendido. Plagio, incluso si viene de un modelo de lenguaje,
          es plagio.
\end{enumerate}
\end{guia}

Se reconoce el uso de dos asistentes de inteligencia artificial
(Claude, de Anthropic, y Codex/ChatGPT, de OpenAI) a lo largo del
proyecto, para tareas de apoyo: auditoría y revisión de código,
implementación técnica puntual bajo pedido y revisión previa a cada
commit, estructuración y redacción de la documentación técnica y
académica, localización de fuentes bibliográficas y estadísticas
oficiales, y generación de borradores de redacción para partes de este
informe.

Las decisiones sobre el alcance del proyecto, la arquitectura, las
fuentes de datos empleadas, los resultados reportados y las conclusiones
son propias. Ningún dato, cifra, cita o resultado incorporado a este
documento se presenta sin verificación contra su fuente primaria o
contra la evidencia técnica del repositorio (código, pruebas
automatizadas, registros de ejecución). La inteligencia artificial no se
utilizó como sustituto de trabajo empírico: no se atribuyen a
entrevistas, usuarios, pruebas, costos o despliegues resultados que no
estén efectivamente documentados.

El registro de trazabilidad con fecha, herramienta, tarea asistida y
verificación aplicada se mantiene en la bitácora de desarrollo del
proyecto.
